\section{沿轨迹的相位为何产生 Bessel 能量梳}
\label{chap:feqo-classical-pinem}

电子沿 \(+z\) 方向以中心速度 \(v_0\) 经过局域近场。光子诱导近场电子显微术（photon-induced near-field electron microscopy，PINEM）中，电子的离散能量边带正来自这样的相互作用。这里固定电场的记法：
\(E_z(z,t)=E_z^{(+)}(z)e^{-\ii\omega t}+E_z^{(+)*}(z)e^{\ii\omega t}\)。
若实验给出的复相量 \(\mathcal E_z\) 满足 \(E_z=\operatorname{Re}[\mathcal E_ze^{-\ii\omega t}]\)，则 \(E_z^{(+)}=\mathcal E_z/2\)。漏掉这个二分之一会把全部 Bessel 函数的自变量错放大一倍。

在 \(\phi=0\) 的规范中 \(E_z=-\partial_tA_z\)，所以正频矢势为 \(A_z^{(+)}=E_z^{(+)}/(\ii\omega)\)。要知道为何能沿一条轨迹积相位，先把第~\ref{chap:feqo-electron-state} 章的正能电子限制为沿 \(+z\) 运动的窄动量波包。令 \(\hat p_z=p_0+\delta\hat p\)，并把中心载波提出为 \(\Psi=e^{\ii(p_0z-E_0t)/\hbar}f(z,t)\)。自由色散给 \(E(p_0+\delta p)=E_0+v_0\delta p+O(\delta p^2/(\gamma_0^3m))\)；在正能分支投影后，Dirac 电流矩阵元 \(c\alpha_z\) 对窄包可用中心群速度 \(v_0\) 代替，故最小耦合项在这个阶次成为 \(-qv_0A_z\)。保留一阶动量偏差时，包络满足
\begin{equation}
\ii\hbar(\partial_t+v_0\partial_z)f(z,t)
=-qv_0A_z(z,t)f(z,t).
\label{eq:feqo-pinem-transport}
\end{equation}
沿特征线 \(z=v_0(t-t_0)\)，左边就是 \(\ii\hbar\dd f/\dd t\)，所以从入射到出射的解是作用势的指数。若包的动量宽度为 \(\delta p\)，这一步要求作用时间 \(T\) 内的自由二次色散相位 \(T(\delta p)^2/(2\gamma_0^3m\hbar)\ll1\)，并要求 \(|\delta p|/(\gamma_0^3mv_0)\ll1\) 使速度矩阵元近似一致。若实际占据到 \(|\ell|\leq N\)，也必须把 \(N\hbar|k_z|\) 计入动量偏离，并检查式~\eqref{eq:feqo-recoil-expansion} 的相位误差；若场强使 \(|qA_z|\) 接近电子动量尺度，或产生显著反射、横向偏转，还需回到原 Dirac 或完整波包方程。现在沿轨迹积分得到
\begin{align*}
U(t_0)&=\exp\!\left[-\frac{\ii}{\hbar}\int V(t)\dd t\right]
=\exp\!\left[\frac{\ii q}{\hbar}\int A_z(z,t_0+z/v_0)\dd z\right]\\
&=\exp\!\left[\beta e^{-\ii\omega t_0}-\beta^*e^{\ii\omega t_0}\right],
\end{align*}
其中无量纲复耦合定义为
\begin{equation}
\beta:=\frac{q}{\hbar\omega}\int_{-\infty}^{\infty}\dd z\,
E_z^{(+)}(z)e^{-\ii\omega z/v_0}.
\label{eq:feqo-beta}
\end{equation}
积分中的 \(qE_z^{(+)}\dd z\) 是能量；指数比较电子到达该位置时的光学相位。若场的空间因子是 \(e^{\ii kz}\) 且只在 \([-L/2,L/2]\) 非零，积分正比于 \(L\sinc[(k-\omega/v_0)L/2]\)：场越长，允许的波数失配越窄。

矩形窗的突然边界不是必需的。取一个可积分的局域近场 \(E_z^{(+)}(z)=E_*e^{-z^2/(2d^2)}e^{\ii kz}\)，其中 \(E_*\) 的单位是 \(\si{\volt\per\metre}\)，\(d\) 是衰减长度。设 \(\Delta k=k-\omega/v_0\)，完成平方并使用高斯积分，得到
\begin{equation}
\beta_G=\frac{qE_*\sqrt{2\pi}d}{\hbar\omega}
e^{-d^2\Delta k^2/2}.
\label{eq:feqo-gaussian-near-field-coupling}
\end{equation}
这一式同时提供量纲和极限检查：\(qE_*d/(\hbar\omega)\) 无量纲；\(\Delta k=0\) 时沿途场相位同向累积；\(|\Delta k|d\gg1\) 时正负贡献相消。可见同一积分耦合同时包含近场幅度、作用长度和速度失配；单一能谱不能把它们各自唯一分开。

再把两个相同的局域作用区放在 \(z=\pm D/2\)，令后一束光相对前一束多一个可控相位 \(\chi\)。变量替换 \(u=z-z_j\) 显示每区的积分多乘 \(e^{\ii\Delta k z_j}\)，因此在区间传播色散相位可忽略时
\[
\beta_{\rm two}=\beta_G
\left(e^{-\ii\Delta kD/2}+e^{\ii\chi}e^{\ii\Delta kD/2}\right),
\qquad
|\beta_{\rm two}|^2=4|\beta_G|^2
\cos^2\!\left(\frac{\Delta kD+\chi}{2}\right).
\]
调 \(\chi\) 可以使两次相位积累相长或相消，即使每一区单独都有非零耦合。若区间传播使不同电子能量分量取得不可忽略的相对相位，两个作用算符之间必须插入传播算符，上式的简单复数相加会失效。

写 \(\beta=|\beta|e^{\ii\phi}\)。由于指数是 \(2\ii|\beta|\sin(\phi-\omega t_0)\)，Jacobi--Anger 恒等式 \(e^{\ii x\sin u}=\sum_{\ell\in\mathbb Z}J_\ell(x)e^{\ii\ell u}\) 给
\begin{equation}
U(t_0)=\sum_{\ell\in\mathbb Z}
J_\ell(2|\beta|)e^{\ii\ell\phi}e^{-\ii\ell\omega t_0},
\qquad
c_\ell=J_\ell(2|\beta|)e^{\ii\ell\phi}.
\label{eq:feqo-pinem-sidebands}
\end{equation}
系数中 \(e^{-\ii\ell\omega t_0}\) 把中心能量 \(E_0\) 移到 \(E_0+\ell\hbar\omega\)；负 \(\ell\) 表示电子失去能量。第 \(\ell\) 条边带的概率是 \(P_\ell=J_\ell^2(2|\beta|)\)。它不含 \(\phi\)，但振幅含有；所以单次能谱无法读场相位，二次相干作用却可以。

回到刚算出的高斯近场，若 \(|\beta_G|\ll1\)，\(J_1(2|\beta_G|)=|\beta_G|+O(|\beta_G|^3)\)，故 \(P_{\pm1}=|\beta_G|^2+O(|\beta_G|^4)\)。双作用区的第一边带则是 \(4|\beta_G|^2\cos^2[(\Delta kD+\chi)/2]+O(|\beta_G|^4)\)。这条从场形状算到可见边带的链条也给出失效条件：\(|\beta_G|\) 不小时必须用完整 Bessel 函数，区间色散不可忽略时必须把传播插在两次作用中。

同一个 \(\beta\) 锁定每个 \(c_\ell\)；边带不能任意独立调节。若 \(2|\beta|\) 扫过 \(J_1\) 的一个零点，\(\ell=1\) 的人口会消失再出现，而总概率仍留在其他通道。这是多路径相干相消，不是光场停止与电子相互作用。

\begin{proposition}[Bessel 边带的归一化与两个矩]
对 \(P_\ell=J_\ell^2(2|\beta|)\)，有
\begin{equation}
\sum_\ell P_\ell=1,\qquad
\sum_\ell\ell P_\ell=0,\qquad
\sum_\ell\ell^2P_\ell=2|\beta|^2.
\label{eq:feqo-bessel-moments}
\end{equation}
\end{proposition}
\begin{proof}
令 \(u(\theta)=e^{\beta e^{-\ii\theta}-\beta^*e^{\ii\theta}}=\sum c_\ell e^{-\ii\ell\theta}\)。指数纯虚，故 \(|u|=1\)。周期 Fourier 正交性给 \(\sum|c_\ell|^2=(2\pi)^{-1}\int_0^{2\pi}|u|^2\dd\theta=1\)。由 \(J_{-\ell}(x)=(-1)^\ell J_\ell(x)\)，正负 \(\ell\) 的概率相等，故一阶矩为零。再将 \(\partial_\theta u=(-\ii\beta e^{-\ii\theta}-\ii\beta^*e^{\ii\theta})u\) 代入 Parseval 等式
\(\sum\ell^2|c_\ell|^2=(2\pi)^{-1}\int|\partial_\theta u|^2\dd\theta\)，交叉项的周期平均为零，余下 \(2|\beta|^2\)。
\end{proof}

这一结果要求电子在交互区主要前向传播，且实际占据边带 \(|\ell|\leq N\) 满足式~\eqref{eq:feqo-recoil-expansion} 下的相位误差远小于一。若横向偏转、反射或反冲可分辨，就需返回包含空间运动或精确反冲的方程；不能把 Bessel 公式解释为任意强场、任意几何的定律。

\begin{exercise}
从 \(J_0(x)=1-x^2/4+O(x^4)\)、\(J_1(x)=x/2+O(x^3)\) 求弱场下 \(P_0,P_{\pm1}\)，核对总概率到 \(|\beta|^2\) 阶为一。
\end{exercise}
